\section{引言}

一位金融分析员问模型：这句话能否被一份财报支持？一位安全工程师问：这段 C
函数是否含有漏洞？一位数学家问：第一个无效步骤出现在哪里？这些场景需要的
是一个可以据以分支的决策，而不是一段必须由人再解读的文字。System One 类
模型直接回应这一需求。给定一个状态和一个带类型的问题，Jev 返回二元
\noul{}、候选 \choice{} 或带概率信息的有序 \score{}，调用程序无需第二次
解释就能据此行动 \citep{typesafe2026,deusser2026jev}。Laya 与 Vela 以开放
权重提供类似接口 \citep{laya2026,vllm2026vela}。不确定性由此成为软件可以
直接读取的一等值。

这个值提示了一条显而易见的控制策略：先问廉价的二元问题，检查返回的概率，
只有当第一个答案看起来不稳时，才把第二次调用花在表述不同的 \choice{} 问题
上。策略容易实现，部署时不需要标签，也符合从业者使用置信度阈值的习惯。但
它依赖一个很少被明说的假设：当前决策所附的概率，能告诉我们另一个决策的
价值。模型可以不确定却仍然正确，也可以很自信却正要从正确的 \noul{} 答案
切换到错误的 \choice{} 答案。路由关心的不是 $P(\text{当前答案正确})$，而是
替代动作的增量效用 $U(\choice\mid x)-U(\noul\mid x)$。

这两个量容易被混为一谈，因为它们是同一标签集上的概率，却由不同的监督信号
学得。RLCD 是 Jev 背后的训练配方，也是 Laya 文档所描述的配方，它用严格
适当评分规则优化所报告分布的诚实性
\citep{typesafe2026,laya2026,gneiting2007scoring}。它教会模型在约七成正确时
说 $0.7$。这一目标从未要求模型预测：换一个 \choice{} 问题是否会得到更好
的答案。因此，一个校准良好的输出头完全可以与一个没有可用信号的路由器
共存。级联继承了这个问题：\citet{jitkrittum2023cascade} 证明，基于置信度
的推迟只有在关于下游误差与上游置信度关系的严格假设下才最优；专门化的下游
误差、标签噪声或分布偏移都可能破坏该规则。

我们测量这一区分，而不是假设它。两个结构化决策模型 Laya 0.4.16 与 Vela 2.0
0.3B 被冻结，并在三个具有不同决策接口的任务上研究推理期路由：CodeXGLUE
缺陷检测 \citep{lu2021codexglue}、FinDVer 金融主张核验
\citep{zhao2024findver} 与 ProcessBench 推理错误定位
\citep{zheng2025processbench}。在二元任务上，动作集为
\{\noul, \choice\}：策略总是先支付 \noul{}，并可能再次支付 \choice{}。在
ProcessBench 上，策略决定是回答 top-1 步骤还是弃权，依据候选概率与
top-1/top-2 间隔。带策略分别用遗传算法（GA）、进化策略（ES）与 NSGA-II
搜索，并在相同数量的唯一候选评估下与固定策略及随机搜索比较。所有输出在
CPU 上生成一次，并从不可变缓存回放，因此每次比较看到的行完全相同。

设计把问题拆成三个。RQ1 问动作集是否含有净互补性：第二个动作修复的行是否
多于它破坏的行？RQ2 问第二个动作的增量效用能否在它执行之前，由当时可用的
信号预测？RQ3 问演化搜索能否把已有信号转化为在留出数据上胜过固定与随机
基线的策略。RQ1 为负会让 RQ3 失去意义；RQ1 为正而 RQ2 为负，则指向表示或
监督问题，而不是优化器问题。两个边界检查——跨任务迁移与按置信度分层的
故障图——检验结论是否只是某个池或某个配置的产物。

结果一致为负，而失败的方式提供了信息。总互补性真实存在：在 CodeXGLUE 与
Laya 上，逐行选择更好动作的 oracle 比更好的固定动作高 23.3 个准确率点，
\choice{} 修复了 14 个 \noul{} 错误。但它同时破坏了 15 个正确答案，因此
训练池上的净修复为 $-1$。留出池并不能稳定这一图景。只有一个 confirm 单元
（FinDVer 与 Laya）的净修复准确率区间为正，但其 macro-F1 区间跨零，且更好
的固定动作在 dev 与 confirm 之间发生反转。效用预测重复了同样的模式：折外
AUROC 在 FinDVer 与 Vela 上达到 0.92，但在 dev 上选出的阈值在 confirm 上
产生负净效用；CodeXGLUE 与 Laya 上最好的 dev 策略退化为对每一行都切换。
搜索也救不了信号：GA 在 60 组配对 macro-F1 比较中没有一次胜过随机搜索；
它对固定切换的 5 次显著胜利全部发生在 FinDVer/Vela 训练池；NSGA-II 的部署
在 6 个单元中显著差于固定基线。校准很差——confirm 的 ECE 从 0.09 到
0.29——但温度缩放让路由结果不变或更差，因为单调重校准无法创造带策略类
原本表达不了的信号。

本文的贡献是把路由价值分解为动作互补性、增量效用可预测性与搜索有效性三个
诊断量，并在冻结的结构化决策模型上测量，而不是从端任务准确率反推。我们
报告一个预注册的负结果：在五个种子、三个池、四个任务—模型单元上使用配对
分组自助法区间，并指出表面胜利只存在于训练池。我们还表明校准与动作价值
预测会分离，并把失败定位在动作集与监督目标，而不是优化器容量。
